\chapter{QCD Renormalization, Asymptotic Freedom, and Deep Inelastic Scattering}
\label{ch:volume-ii-19}

\section{The Running Coupling}

For Yang--Mills theory with gauge group \(SU(N_c)\) and \(N_f\) Dirac fermions
in the fundamental representation, the one-loop beta function has the form
\[
  \mu\frac{\dd g}{\dd\mu}
  =
  -\frac{g^3}{16\pi^2}
  \left(\frac{11}{3}N_c-\frac{2}{3}N_f\right)
  +O(g^5).
\]
For sufficiently small \(N_f\), the coefficient is positive and the coupling
decreases at high energy. This is asymptotic freedom.

The RG-invariant scale is
\[
  \Lambda_{\rm QCD}
  =
  \mu\exp\left(-\frac{8\pi^2}{b_0g^2(\mu)}\right)
  \times(\text{higher-order factors}).
\]
Perturbation theory is reliable at momenta much larger than this scale and
breaks down in the infrared.

\section{Confinement}

The strong-coupling infrared is believed to confine color. Colored fields are
not asymptotic particle states; hadrons are. This is not a minor complication
for scattering theory. It means that LSZ applied to quark and gluon fields
does not define a physical S-matrix. Physical observables must be color
singlet.

\section{Deep Inelastic Scattering}

Deep inelastic scattering probes hadrons at short distances. Current-current
correlators at large momentum transfer admit an operator product expansion:
\[
  J^\mu(x)J^\nu(0)
  \sim
  \sum_a C_a^{\mu\nu}(x,\mu)\mathcal O_a(0,\mu).
\]
The Wilson coefficients are perturbative at short distances, while hadronic
matrix elements of the operators encode long-distance structure. This is the
logic behind parton distributions.

\section{Factorization}

Factorization separates scales without pretending that quarks are free
asymptotic particles. A high-energy observable is expressed as perturbative
short-distance kernels convolved with universal nonperturbative distributions.
The RG evolution of these distributions is controlled by anomalous dimensions
of light-ray or twist operators.
